% !PRINTED 481
\weberpage{481}
%<<CONT
Denn die als unabhängige Variable betrachteten Co\"efficienten
$a_1, a_2 \ldots a_m$ können auch dargestellt werden als die symmetrischen
Grundfunctionen der Wurzeln $\alpha, \alpha_1, \alpha_2 \ldots \alpha_{m-1}$.
Ist dann $g(t)$ ein rational durch die $a$ ausgedrückter, irreducibler Factor
von $G(t)$, der für $t = \varrho$ verschwindet, so erhält man, wenn man $a_i$
und $\varrho$ durch die $\alpha$ darstellt, in $g(\varrho) = 0$ eine identische
Gleichung. In dieser können aber die $\alpha$ beliebig permutirt werden,
wodurch sich die $a$ nicht ändern, während $\varrho$ in jede der Wurzeln
$\varrho, \varrho', \varrho'' \ldots$ von $G(t)$ übergehen kann. Es
verschwindet also $g(t)$ für alle Wurzeln von $G(t)$ und muss folglich mit
$G(t)$ übereinstimmen\footnote{Évariste \gesperrt{Galois} ist im Jahre 1832,
kaum 20jährig, im Duell gefallen. Die erste Andeutung über die Theorie, die
heute seinen Namen trägt, findet sich in einer 1830 im „Bulletin des sciences
mathém.“ von \gesperrt{Férussac} erschienenen Abhandlung „Analyse d'un mémoire sur la
résolution algébrique des équations“. Ausführlichere Mittheilungen enthält
der berühmte, am Vorabend seines Todes geschriebene Brief an \gesperrt{Auguste
Chevalier}, der in der „Revue encyclopédique“ vom September 1832
veröffentlicht ist. Erst im Jahre 1846 hat \gesperrt{Liouville} die sämmtlichen schon
veröffentlichten Arbeiten von \gesperrt{Galois} nebst einigen Untersuchungen aus dem
Nachlass, darunter die wichtigste „Mém. sur les conditions de résolubilité
des équations par radicaux“ in Bd. XI seines Journals abdrucken lassen. Eine
deutsche Ausgabe von \gesperrt{Maser} ist 1889 erschienen (Berlin, bei Springer).
Interessante biographische Mittheilungen finden sich in der „Revue
encyclopédique“ vom September 1832, die \gesperrt{Liouville} nicht mit abgedruckt hat
und die auch in der deutschen Ausgabe fehlen. Von weiteren Schriften, die zum
Verständniss oder zur Weiterbildung der \gesperrt{Galois'schen} Theorie beigetragen
haben, seien noch erwähnt: J. A. \gesperrt{Serret} („Cours d'algèbre supérieure“,
II. Ausgabe, 1854, IV. 1879); \gesperrt{Betti} (Annali di scienze fisiche e matematiche
[1853]); C. \gesperrt{Jordan}, „Traité des substitutions“ (1870). E. \gesperrt{Netto},
„Substitutionentheorie“ (Leipzig 1882).}.

\begin{center}
\S.\,150.

\gesperrt{Transitive und intransitive Gruppen.}
\end{center}

Aus der \gesperrt{Galois'schen} Gruppe $P$ können wir ein sehr
einfaches Kennzeichen dafür herleiten, ob die Gleichung $F(x) = 0$ reducibel
oder irreducibel ist.

Nehmen wir an, es sei $F(x)$ vom Grade $m$ reducibel und $f(x)$ ein Factor
von $F(x)$ in $\Omega$ vom Grade $n$ und $n < m$. Es mögen die Wurzeln von
$f(x)$
%CONT>>
